% ===== 12. Наступні кроки =====
\begin{frame}{Наступні кроки}
  \begin{enumerate}
    \setlength{\itemsep}{6pt}
    \item \textbf{Обсяг MVP:} відео GoPro, одна відкрита траса, 2–3 поради
    \item \textbf{Технології:} перевірити gopro-telemetry
    \item \textbf{Методика демо:} що автоматизуємо — межі, зсув GPS
    \item \textbf{Супутникові знімки:} перевірити умови використання
    \item \textbf{Конкуренти:} встановити Tenth Finder і Kart Track
  \end{enumerate}

  \note{[Доповідач 4] Що далі. Спочатку визначаємо обсяг MVP: відео GoPro, одна відкрита траса, поділ на кола, втрати часу по поворотах і дві-три поради. Потім обираємо технології і першою перевіряємо бібліотеку gopro-telemetry. Також розбираємо методику нашого демо і вирішуємо, що з неї робимо автоматично: обведення меж траси за супутником і корекцію зсуву GPS. Окремо треба перевірити умови використання супутникових знімків. І ще до розробки встановлюємо Tenth Finder і Kart Track, щоб зрозуміти, яке обладнання їм потрібне і чим ми від них відрізняємося.}
\end{frame}

% ===== 13. Дякуємо =====
\begin{frame}[standout]
  Дякуємо за увагу!\\[0.3em]
  Запитання?

  \vspace{1.5em}
  \raggedright\footnotesize\mdseries
  Повний SWOT і джерела:\\
  \texttt{docs/stages/stage-1/lap-coach/SWOT.md}\\
  \texttt{docs/stages/stage-1/smart-grocery/SWOT.md}\\
  Рішення, критерії, ризики: \texttt{docs/stages/stage-1/Decision.md}\\
  Макет розбору заїзду: \texttt{docs/stages/stage-1/lap-coach/Karting-session.html}

  \note{[Доповідач 4] Дякуємо за увагу! Коротко ще раз: ми розглянули дві теми, описали їх, зробили SWOT-аналіз, порівняли за п'ятьма критеріями й обрали цифровий тренер для автоспорту. Він виграє за конкурентною перевагою і меншою залежністю від однієї компанії. Внизу слайда — шляхи до файлів із повним SWOT-аналізом, джерелами і документом з рішенням: там є всі цифри, які ми сьогодні називали, з посиланнями. Якщо цікаво, можемо відкрити макет розбору заїзду в браузері. Будемо раді вашим запитанням.}
\end{frame}
